\documentclass[10pt]{article}

\usepackage[a4paper,margin=0.65in]{geometry}
\usepackage{amsmath}
\usepackage{amssymb}
\usepackage{booktabs}
\usepackage{array}
\usepackage{enumitem}
\usepackage{multicol}
\usepackage{xcolor}
\usepackage{hyperref}

\hypersetup{
    colorlinks=true,
    linkcolor=black,
    urlcolor=blue
}

\setlength{\parindent}{0pt}
\setlength{\parskip}{0.25em}
\setlist[itemize]{leftmargin=1.2em,itemsep=0.12em,topsep=0.15em}
\setlist[enumerate]{leftmargin=1.4em,itemsep=0.12em,topsep=0.15em}
\renewcommand{\arraystretch}{1.12}

\DeclareMathOperator{\Var}{Var}
\DeclareMathOperator{\Cov}{Cov}
\DeclareMathOperator{\E}{E}
\newcommand{\dd}{\Delta}
\newcommand{\pnl}{\mathrm{P\&L}}

\definecolor{accent}{HTML}{204A87}
\definecolor{soft}{HTML}{EEF4FB}

\newcommand{\formula}[1]{\[
    \boxed{\displaystyle #1}
\]}
\newcommand{\note}[1]{\textcolor{accent}{\textbf{Check:}} #1}

\begin{document}

\begin{center}
    {\Large \textbf{Derivatives and Alternative Investments}}\\
    \vspace{0.2em}
    {\large Exam Formula-Check Notes}\\
    \vspace{0.2em}
    \small Coverage: Lecture 1, Lecture 2 sections 1--3 only, Lecture 4. Lecture 3 is excluded.
\end{center}

\vspace{0.4em}

\begin{multicols}{2}

\section*{Notation}
\begin{tabular}{@{}ll@{}}
\toprule
Symbol & Meaning \\
\midrule
$S_t$ & Spot price at time $t$ \\
$S_T$ & Spot price at maturity/expiration \\
$F_t$ & Forward or futures price at time $t$ \\
$F_0$ & Initial forward/futures price \\
$K$ & Delivery price or option strike \\
$T$ & Time to maturity, in years \\
$r$ & Continuously compounded risk-free rate \\
$q$ & Continuous income/dividend yield \\
$I$ & Present value of known cash income \\
$U$ & Present value of known storage cost \\
$u$ & Proportional storage-cost rate \\
$y$ & Convenience yield \\
$c,p$ & European call and put values \\
$C,P$ & American call and put values \\
$N$ & Number of futures contracts \\
$Q$ & Contract size \\
$h$ & Hedge ratio \\
\bottomrule
\end{tabular}

\note{Most pricing formulas below assume continuous compounding unless stated otherwise.}

\section*{1. Payoffs and Contract Values}

\subsection*{Derivative Basics}
A derivative payoff can be written generally as
\formula{\text{payoff at }T=g(X_T)}
where $X_T$ is the underlying variable at maturity.

\textbf{Payoff} is the terminal cash flow. \textbf{Profit} is payoff minus initial cost and financing. \textbf{Value} is the current market value of the contract.

\subsection*{Forwards}
Long forward: agree to buy at $K$.
\formula{\text{long forward payoff}=S_T-K}

Short forward: agree to sell at $K$.
\formula{\text{short forward payoff}=K-S_T}

For a non-income-paying investment asset:
\formula{F_0=S_0e^{rT}}

With simple annual interest, use the matching simple-compounding version, e.g.
\[
    F_0=S_0(1+rT).
\]

\note{If quoted $F_0>S_0e^{rT}$, cash-and-carry: borrow, buy asset, short forward. If quoted $F_0<S_0e^{rT}$, reverse cash-and-carry: short asset, invest proceeds, long forward.}

\subsection*{Futures Daily Settlement}
For $N$ contracts and contract size $Q$:
\formula{\pnl_t^{\text{long}}=NQ(F_t-F_{t-1})}
\formula{\pnl_t^{\text{short}}=NQ(F_{t-1}-F_t)}

\note{Margin is collateral, not the price of the futures contract. Daily settlement reduces credit risk but can create liquidity risk.}

\subsection*{Options}
Call option payoff:
\formula{\max(S_T-K,0)=(S_T-K)^+}

Put option payoff:
\formula{\max(K-S_T,0)=(K-S_T)^+}

Profit, ignoring interest in payoff diagrams:
\[
    \text{long option profit}=\text{payoff}-\text{premium}.
\]

American exercise weakly increases value:
\formula{C\geq c,\qquad P\geq p}

\section*{2. Hedging With Futures}

\subsection*{Choosing Hedge Direction}
\begin{itemize}
    \item \textbf{Short hedge:} use when the firm owns, produces, or expects to receive an asset it will later sell.
    \item \textbf{Long hedge:} use when the firm expects to buy an asset later and wants protection against rising prices.
    \item \textbf{Cross hedge:} use a related futures contract when no liquid futures contract exists on the exact exposure.
\end{itemize}

Simple contract count when hedge ratio is one:
\formula{N=\frac{\text{quantity to hedge}}{\text{contract size}}}

\subsection*{Basis Risk}
Basis is spot minus futures:
\formula{b_t=S_t-F_t}

For a short hedge opened at $t=1$ and closed at $t=2$:
\[
    \text{futures gain}=F_1-F_2,
\]
\formula{\text{effective sale price}=S_2+(F_1-F_2)=F_1+b_2}

For a long hedge opened at $t=1$ and closed at $t=2$:
\[
    \text{futures gain}=F_2-F_1,
\]
\formula{\text{effective purchase price}=S_2-(F_2-F_1)=F_1+b_2}

\note{If $b_2=0$, the hedge locks in $F_1$. Otherwise, $b_2$ is the remaining basis risk.}

Cross-hedge basis decomposition:
\formula{S_2-F_2=(S_2-S_2^*)+(S_2^*-F_2)}
The first term is mismatch between physical assets; the second is the usual basis for the futures underlying.

\subsection*{Minimum-Variance Hedge Ratio}
For a short hedge using hedge ratio $h$:
\[
    \Delta V=\Delta S-h\Delta F.
\]

Residual variance:
\[
    \Var(\Delta V)=\sigma_S^2+h^2\sigma_F^2-2h\Cov(\Delta S,\Delta F).
\]

Optimal hedge ratio:
\formula{h^*=\frac{\Cov(\Delta S,\Delta F)}{\Var(\Delta F)}
=\rho\frac{\sigma_S}{\sigma_F}}

Contract count:
\formula{N^*=h^*\frac{Q_A}{Q_F}
\qquad\text{or}\qquad
N^*=h^*\frac{V_A}{V_F}}

\note{If the regression uses price changes, use physical quantities. If it uses percentage returns, use values.}

\section*{3. Forward and Futures Pricing}

\subsection*{Investment Assets}
Non-income-paying asset:
\formula{F_0=S_0e^{rT}}

Known cash income with present value $I$:
\formula{F_0=(S_0-I)e^{rT}}

Known continuous yield $q$:
\formula{F_0=S_0e^{(r-q)T}}

Value of an existing long forward with delivery price $K$:
\formula{f=(F_0-K)e^{-rT}}

For an existing short forward:
\[
    f_{\text{short}}=(K-F_0)e^{-rT}.
\]

\subsection*{Commodity Pricing}
Fixed storage cost with present value $U$:
\formula{F_0=(S_0+U)e^{rT}}

Proportional storage cost $u$:
\formula{F_0=S_0e^{(r+u)T}}

With proportional storage cost $u$ and convenience yield $y$:
\formula{F_0=S_0e^{(r+u-y)T}}

For consumption commodities, arbitrage may give only an upper bound:
\formula{F_0\leq (S_0+U)e^{rT}}

\textbf{Contango:} futures prices rise with maturity. \textbf{Backwardation:} futures prices fall with maturity.

\subsection*{Futures Prices and Expected Spot Prices}
The futures price is not just a forecast:
\formula{F_0=\E(S_T)e^{(r-k)T}}
where $k$ is the required expected return on the underlying asset.

\begin{center}
\begin{tabular}{@{}lll@{}}
\toprule
Systematic risk & Relation & Implication \\
\midrule
None & $k=r$ & $F_0=\E(S_T)$ \\
Positive & $k>r$ & $F_0<\E(S_T)$ \\
Negative & $k<r$ & $F_0>\E(S_T)$ \\
\bottomrule
\end{tabular}
\end{center}

\note{Do not confuse backwardation/contango as curve shapes with normal backwardation/contango as statements about expected future spot prices.}

\section*{4. Options and Moneyness}

\subsection*{Four Basic Positions}
\begin{center}
{\footnotesize
\begin{tabular}{@{}lll@{}}
\toprule
Position & Payoff & Exposure \\
\midrule
Long call & $\max(S_T-K,0)$ & Upside \\
Short call & $-\max(S_T-K,0)$ & Upside loss \\
Long put & $\max(K-S_T,0)$ & Downside protection \\
Short put & $-\max(K-S_T,0)$ & Downside loss \\
\bottomrule
\end{tabular}
}
\end{center}

\subsection*{Break-Even Points}
Ignoring interest:
\[
    \text{long call break-even}=K+\text{premium},
\]
\[
    \text{long put break-even}=K-\text{premium}.
\]

\subsection*{Moneyness}
\begin{center}
\begin{tabular}{@{}lll@{}}
\toprule
 & Call & Put \\
\midrule
In the money & $S>K$ & $S<K$ \\
At the money & $S=K$ & $S=K$ \\
Out of the money & $S<K$ & $S>K$ \\
\bottomrule
\end{tabular}
\end{center}

Intrinsic value:
\formula{\text{call intrinsic value}=\max(S-K,0)}
\formula{\text{put intrinsic value}=\max(K-S,0)}

Time value:
\formula{\text{time value}=\text{option price}-\text{intrinsic value}}

A standard U.S. exchange-traded stock option contract usually covers 100 shares, and quoted option prices are usually per share.

\section*{5. Option Price Restrictions}

\subsection*{Inputs to Stock Option Prices}
\begin{center}
{\footnotesize
\begin{tabular}{@{}lll@{}}
\toprule
Input increases & Call & Put \\
\midrule
$S_0$ & Increases & Decreases \\
$K$ & Decreases & Increases \\
$\sigma$ & Increases & Increases \\
$r$ & Usually increases & Usually decreases \\
Dividends & Decreases & Increases \\
\bottomrule
\end{tabular}
}
\end{center}

\subsection*{Upper Bounds}
\formula{c\leq S_0,\qquad C\leq S_0}
\formula{p\leq K,\qquad P\leq K}
For a European put:
\formula{p\leq Ke^{-rT}}

\subsection*{Lower Bounds}
For European options on a non-dividend-paying stock:
\formula{c\geq \max(S_0-Ke^{-rT},0)}
\formula{p\geq \max(Ke^{-rT}-S_0,0)}

\section*{6. Put-Call Parity}

For European options on a non-dividend-paying stock:
\formula{c+Ke^{-rT}=p+S_0}

Equivalent form:
\formula{c-p=S_0-Ke^{-rT}}

Fiduciary call:
\[
    \text{call}+\text{cash }Ke^{-rT}.
\]
Protective put:
\[
    \text{put}+\text{stock}.
\]
Both pay $\max(S_T,K)$ at maturity.

\subsection*{Parity Arbitrage Rule}
Compare the two packages:
\[
    \underbrace{c+Ke^{-rT}}_{\text{fiduciary call}}
    \quad\text{and}\quad
    \underbrace{p+S_0}_{\text{protective put}}.
\]

\begin{center}
\begin{tabular}{@{}p{0.36\linewidth}p{0.54\linewidth}@{}}
\toprule
Mispricing & Arbitrage direction \\
\midrule
$c+Ke^{-rT}>p+S_0$ &
Sell fiduciary call; buy protective put \\
$c+Ke^{-rT}<p+S_0$ &
Buy fiduciary call; sell protective put \\
\bottomrule
\end{tabular}
\end{center}

\note{Buy the cheap package and sell the expensive package. The maturity payoffs cancel; the mispricing is cash today.}

\section*{7. Binomial Trees}

\subsection*{One-Step Setup}
Stock movement:
\formula{S_u=S_0u,\qquad S_d=S_0d,\qquad u>1,\quad d<1}

Terminal option values are $f_u$ and $f_d$.

\subsection*{Delta and Replication}
Make $\Delta$ shares minus one option riskless:
\[
    S_0u\Delta-f_u=S_0d\Delta-f_d.
\]

Thus:
\formula{\Delta=\frac{f_u-f_d}{S_0u-S_0d}
=\frac{f_u-f_d}{S_0(u-d)}
=\frac{f_u-f_d}{S_u-S_d}}

\note{Delta is local to the current node and time step. In a multi-step tree it changes across nodes and must be rebalanced.}

\subsection*{Risk-Neutral Probability and Price}
One-step option price:
\formula{f=e^{-rT}\left[p f_u+(1-p)f_d\right]}

Risk-neutral probability:
\formula{p=\frac{e^{rT}-d}{u-d}}

Equivalent condition defining $p$:
\formula{pS_0u+(1-p)S_0d=S_0e^{rT}}

\note{$p$ is not generally the real-world probability of an up move. It is the pricing probability that makes the stock earn the risk-free rate.}

\subsection*{Multi-Step Backward Induction}
If option life is $T$ and there are $n$ steps:
\formula{\Delta t=\frac{T}{n}}

At each node:
\formula{f=e^{-r\Delta t}\left[p f_u+(1-p)f_d\right]}

For two steps:
\[
    f_u=e^{-r\Delta t}\left[p f_{uu}+(1-p)f_{ud}\right],
\]
\[
    f_d=e^{-r\Delta t}\left[p f_{ud}+(1-p)f_{dd}\right],
\]
\[
    f=e^{-r\Delta t}\left[p f_u+(1-p)f_d\right].
\]

Two-step terminal stock prices:
\formula{S_0u^2,\qquad S_0ud,\qquad S_0d^2}

\subsection*{Choosing Tree Parameters}
For volatility $\sigma$ and time step $\Delta t$:
\formula{u=e^{\sigma\sqrt{\Delta t}},\qquad
d=e^{-\sigma\sqrt{\Delta t}}=\frac{1}{u}}

Then:
\formula{p=\frac{e^{r\Delta t}-d}{u-d}}

\note{Given $S_0,u,d,r$ and the payoff, the option price is pinned down by arbitrage. The stock's real-world expected return does not enter.}

\end{multicols}

\end{document}
